\section{Methods}
\label{sec:methods}

\subsection{Encoder and head}

Every arm uses the Zoobot greyscale encoder~\citep{zoobot}, a ConvNeXt-Nano loaded from the
Hugging Face hub. The Zoobot 2.0 encoders start from ImageNet-12k weights and are trained on
about \zoobotPretrainImages{} Galaxy Zoo--labelled images from ground-based and HST surveys,
none of them Euclid~\citep{walmsley2024scaling}; the greyscale encoder we use was released
later, and its training is not described separately. The encoder returns a 640-dimensional feature
vector. Every arm passes it through a fixed per-dimension standardisation, fitted on cached
source features, before the head. On top sits one head, identical across arms: $\mathrm{LayerNorm}(640)$ followed by a
linear map to a single logit, 1921 trainable parameters. The encoder contains no batch
normalisation and no active dropout, so train and evaluation modes compute the same features.
Appendix~\ref{app:repro} records the software versions.

\subsection{Frozen and unfrozen regimes}

Arms run either \emph{frozen}, with only the head trained, or \emph{unfrozen}, with the top
three timm stages of the encoder trained as well. Because ConvNeXt-Nano is top-heavy, that is
14\,835\,520 trainable encoder parameters, 99.2\% of the encoder, so the unfrozen regime is
close to full fine-tuning; the benchmark fine-tunes the last three blocks of the same
backbone~\citep{lines2025simstosky}, which we implement as these three stages. Unfrozen arms
start from a converged head trained with the encoder frozen and then train the encoder, the order that avoids distorting pretrained features out of
distribution~\citep{kumar2022lpft}. Each unfrozen run is gated at step zero against its head's
recorded value on that fold, which is the check that the run used the right fold.
BNM's pre-registration required both halves of its step-zero gate to match exactly; its runner
applied one positive and $\bnmGateAurocTol{}$ in AUROC, as for DANN. TPR matched exactly on
every fold and AUROC within $\bnmGateAurocBound{}$. All arms
use batch size \rwBatch{} and AdamW. The unfrozen CORAL, DANN, BNM and $\lambda=0$ arms use
the encoder's published fine-tuning defaults (learning rate $\rwOptLr{}$, layer-wise decay
\rwOptDecay{}, weight decay \rwOptWd{}). The unfrozen FreeMatch arm keeps the settings of its
frozen-comparison run (head learning rate $\rwFmHeadLr{}$, encoder learning rate
$\rwFmEncLr{}$, weight decay $\rwFmWd{}$), so that its frozen-to-unfrozen comparison changes
only the encoder setting. Baseline A, which trains only the head, uses learning rate
$\rwBaseLr{}$ and weight decay $\rwBaseWd{}$. We tuned no optimiser setting for any arm, apart
from one comparison: for the CORAL, DANN, BNM and $\lambda = 0$
arms, we compared the encoder's published defaults with one alternative on the development
fold, using the source-only arm, and kept the defaults, which did not score higher. Every step includes a labelled source batch of
\batchSourceRows{} rows, balanced 1:1, next to \batchTargetRows{} target rows in every arm that
loads target images.

\subsection{The arms}
\label{sec:arms}

\paragraph{Baseline A.} Supervised training on the labelled source with the encoder frozen and
no target term. It sees no target image, so one model serves all five folds.

\paragraph{The $\lambda=0$ control.} The DANN arm with its adaptation weight $\lambda$ set to
zero. Target images pass through the discriminator, but the gradient-reversal layer returns
zero gradient, so no target-derived gradient reaches the encoder or the head. The
discriminator itself still trains, and in half precision an overflow in its gradients skips the
whole step; the control skipped \lamSkipFoldTwo{} and DANN \dannSkipFoldTwo{} of 4000 steps on
fold 2, so the control's update schedule is not wholly independent of the target images. DANN and this
control share a head, a budget, a seed and a step grid and differ, in every training-relevant
respect, only in $\lambda$. DANN's confirmatory folds ran with abort conditions armed, while the
control's folds only logged them; no arming condition fired on DANN, and every DANN fold ran its
full 4000 steps. This makes DANN $-$ $\lambda=0$ the only fully controlled pair in this paper. Most other contrasts
also carry one or both of two differences that pairing by fold does not remove. The unfrozen arms start
from three different heads: CORAL from its frozen-comparison heads, FreeMatch from its own,
and DANN, BNM and the control from Baseline A's. Their tuning effort also differs. CORAL's
alignment weight came from a three-point sweep that read only loss magnitudes and feature
statistics. BNM's weight came from a development-fold sweep over a grid set by gradient norms
at step zero; the sweep read the nuclear norm and the predicted-positive fraction, never TPR.
FreeMatch's local-threshold variant was chosen over an alternative by a development-fold
screen of the metric, and DANN carried no tuning. BNM's runner numbers epochs from zero and
DANN's and the control's from one, so on each fold its batches are a different shuffle of the
same source data, and BNM $-$ $\lambda=0$ is not matched in batch order.

\paragraph{BNM.} Batch nuclear-norm maximisation~\citep{cui2020bnm} acts on target predictions
only: no domain comparison, pseudo-label, threshold or discriminator, and one weight,
$\gamma = 0.00787$. Its loss is the nuclear norm divided by $\min(B, C) = \bnmNucDenom{}$, as
in \citeauthor{cui2020bnm}'s code, so $\gamma$ refers to that normalisation. Maximising the nuclear norm of the batch prediction matrix raises both
confidence and prediction diversity~\citep{cui2020bnm}. Before any run, we checked,
analytically and on cached target features, that the objective separates a balanced confident batch from a collapsed one at
our two-class setting.

\paragraph{FreeMatch.} A consistency objective between weakly and strongly augmented target
views, with self-adaptive confidence thresholds and a fairness term~\citep{wang2022freematch}.
We kept FreeMatch's threshold momentum (\fmSatMomentum{}) and fairness weight
(\fmSafWeight{}), but our fairness term compares the mean confident prediction with a fixed
uniform target rather than FreeMatch's running class estimate, and our consistency weight
follows AdaMatch's cosine ramp~\citep{berthelot2022adamatch} over the first half of training,
with no consistency term for the first \fmWarmupPct{}\% of steps; the unfrozen stage starts
after that ramp, so it trains at full weight throughout. Its strong view is the augmentation of
Section~\ref{sec:sim-pipeline}.

\paragraph{DANN.} A domain discriminator behind a gradient-reversal
layer~\citep{ganin2016dann}, mapping $640 \rightarrow 1024 \rightarrow 1024 \rightarrow 1$ with
layer normalisation, ReLU and dropout after each hidden layer; the dropout rate of 0.2 is that
of~\citet{belfiore2024}, and the placement is our choice. The discriminator of \citet{ganin2016dann},
of the same widths, has no normalisation layer; we chose layer rather than batch normalisation so
the discriminator cannot read the domain from batch statistics. It shares the AdamW optimiser,
at learning rate $\dannDiscLr{}$ and weight decay \dannDiscWd{}, without the learning-rate
annealing of \citet{ganin2016dann}. The adaptation weight follows $\lambda = \tanh(5s/2000)$ in the
step $s$ up to step 2000, after which it stays at $\tanh(5)$: their schedule, with training
progress replaced by a fixed horizon of 2000 steps.

\paragraph{CORAL.} We ran the DomainBed variant of correlation
alignment~\citep{gulrajani2021domainbed}, which aligns feature means and covariances rather
than covariances alone~\citep{sun2016coral}. Our head emits a single logit, so we compute the
loss one layer earlier, on the 640-dimensional layer-normalised features, with weight
$\gamma = 16$, after shrinking each covariance \coralShrinkPct{}\% toward a scaled identity; as
in DomainBed, both the mean and covariance terms are averaged over their entries rather than
scaled by \citeauthor{sun2016coral}'s $1/(4d^2)$. With the encoder changing every step, we disabled the running average of the
target covariance.

\subsection{Where labels enter}
\label{sec:labels}

\paragraph{The balanced design.} No real-lens label is ever used as a training target, but
labels shape training in two ways, and we use them a third time for evaluation. First, they
build each fold's target pool: the
$\rvTargetLensPool{}$ graded lenses of the other four folds, without their labels, plus an
equal number of random galaxies. Second, they balance every target batch: our sampler draws
half of each batch from the lens part of the pool and half from the galaxy part, using the
catalogue labels to tell them apart. The DANN, BNM and $\lambda=0$ runs also hold out a
$\rvHoldoutN{}$-image slice of the target pool, chosen half lenses and half galaxies in the
same way, and FreeMatch rebuilds its confidence thresholds from the same balanced loader.
Third, the labels of the held-out fold score the test. Baseline A loads no target images, and
the $\lambda=0$ control passes no target gradient to the encoder or the head, so neither is
trained on the target labels; the control's only dependence on them is through the few
half-precision steps it skips (Section~\ref{sec:arms}). Every adaptation arm does depend on
them. Because the batches are balanced, the class proportion itself reaches every
adaptation objective (Section~\ref{sec:limitations}).

\paragraph{The diluted design.} In our second experiment (Section~\ref{sec:dilute-methods}) the
labels still choose the pool: which \dlPoolLens{} lenses enter it, and the exact count of
\dlRatio{} galaxies per lens (Section~\ref{sec:splits}). Target batches are drawn uniformly at
random from the pool, without the labels, so a batch of \batchTargetRows{} holds about
\dlBatchLens{} lenses on average (about \dlBatchLensTrain{} for DANN and BNM, whose holdout is
removed), and the number varies from batch to batch. The DANN and BNM
runs keep the \dlHoldout{}-image holdout of the balanced design, which is half lenses by
construction, so their training pool is \dlTrainLens{}:\dlTrainGal{}
(\dlTrainLensFour{}:\dlTrainGalFour{} on fold 4) and their batches average a lens fraction of
$\dlTrainFracThree{}$ rather than $\dlPoolFracThree{}$. FreeMatch rebuilds its threshold state
from the uniform sampler.

\paragraph{Why the main design is balanced.} We chose a balanced target because adaptation
failed without one. In early runs with a frozen
encoder and a target far less balanced than the source, DANN and an importance-weighted
variant, IW-CDAN~\citep{zhao2020gls}, both collapsed. On a morphology test (grade-X objects
against random galaxies), alignment lowered the trained head's AUROC from $\rvEarlyNoDa{}$ to
$\rvEarlyDann{}$ for DANN and $\rvEarlyIw{}$ for IW-CDAN, while a linear probe on the same
frozen features reached $\rvEarlyProbe{}$. The information was still in the features, but
alignment removed the head's access to it. Those runs used an earlier pipeline and this proxy
test, and their run artifacts are not in our repository, so we report them as the reason for
the design, not as a result. We then tested part of the design's main weakness
(Section~\ref{sec:results-dilute}): with the encoder unfrozen and a 1:10 target whose batches are
drawn without the labels, DANN keeps a detectable gain over the control, although the comparison
with the balanced target cannot rule out that it lost much of that gain. The early runs differed from these in their frozen encoder,
their pipeline and the degree of imbalance, so the two observations do not contradict each
other.

\subsection{The budget rule}
\label{sec:budget}

We fixed each unfrozen arm's training budget on the development fold, fold 0, and read every arm at its
last checkpoint, so the budget is the model selector. Fold 0 enters no confirmatory number.
BNM's budget came from a rule declared before its development run, and Baseline A's was fixed
before any results; the others were chosen with the development curve in view.
Appendix~\ref{sec:hashes} gives each budget, its basis and its timing. The development fold is
scored against the same shared test negatives as the confirmatory folds, so the budgets we chose
with its curve in view saw those negatives, though none of the confirmatory positives.

\subsection{Pre-registration}
\label{sec:prereg}

We fixed each arm's budget and read-off step before its own confirmatory folds ran, and
hash-protected them then, with four exceptions that rest on our run logs alone
(Appendix~\ref{sec:hashes}). The $\lambda=0$ control was a later addition. On 2026-09-01 we had
decided to run it on the development fold only; on 2026-09-06, after the FreeMatch, CORAL and
DANN confirmatory folds had run, we extended it to folds 1--4, because those arms were unfrozen
and the baseline they were being compared with was not, so their gains could not be attributed
to adaptation. Its budget, read-off rule and decision branches, and the four paired contrasts
against it, were declared and hash-protected before any of its confirmatory folds ran, so the
registration fixes how the control was run and read; the choice of contrasts was made with the
adaptation arms' results known. Appendix~\ref{sec:hashes} lists each file, its digest and its
timing, and the four exceptions.

\subsection{The diluted-target experiment}
\label{sec:dilute-methods}

The second experiment asks whether the gain over the control needs the balanced target. It
changes two things: the target pool, diluted to one lens per \dlRatio{} galaxies
(Section~\ref{sec:splits}), and the sampler, which draws target batches uniformly, without the
labels (Section~\ref{sec:labels}). Nothing else changes: the heads, optimisers, adaptation
weights, seeds, budgets, read-off steps and step-zero gate rules, apart from two tolerances,
are those of the main design, and we retuned nothing.

CORAL and FreeMatch start their unfrozen runs from a head trained with the encoder frozen, so we
reran those frozen stages on the diluted target, with the new pool and sampler and no other
change. Their final heads initialise the unfrozen CORAL and FreeMatch arms, and each unfrozen
run's step-zero gate is taken against its new head. FreeMatch's fold 1, whose balanced run was
gated at an exact tolerance, now uses the one-positive tolerance of the other folds. CORAL's
gate, which in the main design only warned between one positive and \coralGateWarnTpr{} in TPR
or between $\coralGateAurocWarn{}$ and \coralGateAurocAbort{} in AUROC,
now stops a run outside one positive or $\dlGateAurocTol{}$ in AUROC, as for the other arms. DANN and BNM
start from Baseline A's head, as before. We reused the $\lambda=0$ control without rerunning it:
it passes no target gradient to the encoder or the head, so the target pool and sampler enter
what it learns only through the few half-precision steps it skips (Section~\ref{sec:arms}). We
checked that the new sampler shares no random stream with the source sampler or the
augmentation: replaying the samplers alone over the first \dlRngSteps{} training steps of every
arm and fold (for FreeMatch, after its threshold rebuild and warm-up), the source batches under
the new sampler were bit-identical to those under the balanced sampler.

Dilution changes three things at once. It lowers the target's class proportion from 1:1 to
1:\dlRatio{}; it cuts the number of real lenses in a target batch from \dlBatchLensBal{} to
about \dlBatchLens{} (about \dlBatchLensTrain{} for DANN and BNM, whose holdout is removed); and for a fixed number of steps each target lens is seen about
\dlExposureRatio{} times less often. The galaxy side of the target also becomes about
\dlRatio{} times more varied. A change in the gain therefore cannot be attributed to the
proportion alone. The experiment tests whether the gain survives a target far from balanced, as
a first step toward survey prevalence; it does not separate adapting to real images from the
class proportion.

We pre-registered two contrasts per arm, each a paired $t$-interval over folds 1--4 on the
cleaned pool. The primary, $G$, is the diluted run minus the $\lambda=0$ control, with the
control read at 4000 as in the main design's contrasts against the control; for CORAL and FreeMatch,
the same contrast with the control read at the arm's own step is a pre-declared sensitivity
readout. The secondary, $D$, is the balanced run minus the diluted run at the same read-off
step. It carries no verdict, and it includes the noise of a changed batch order. Each arm's
branch follows from $G$ and $D$:
\begin{itemize}
\item \emph{survives}: $G$'s interval excludes zero, and $D$'s includes zero or is negative;
\item \emph{shrinks}: $G$'s interval excludes zero, and $D$'s excludes zero and is positive;
\item \emph{collapses}: $G$'s interval includes zero.
\end{itemize}
Each branch has a fixed sentence, which we report with its numbers filled in
(Section~\ref{sec:results-dilute}). Under \emph{survives}, if $D$'s upper bound exceeds half of
the arm's balanced gain over the control, the sentence must add that the interval admits losing
up to that bound. \emph{Collapses} means only that we cannot detect a gain.

The pre-registration has three parts (Appendix~\ref{sec:hashes}). The first, recorded before
any run, fixes the pool, the sampler, the arms, the contrasts, the branches and our predictions.
The second, recorded after step-zero measurements of at most \dlPartTwoMaxStepsWord{}
training steps per fold, but before any confirmatory run, fixes the step-zero thresholds and which inherited
stop conditions abort a run and which only log. The third was recorded after the first attempt
of CORAL's unfrozen run on fold 1 stopped at its first step (Section~\ref{sec:limitations}), and
before any confirmatory CORAL or FreeMatch unfrozen run. It adds a rule that withholds the verdict from an arm whose run applies fewer than half its
budget's optimiser updates; we set it after a label-free \dlCoralCheckSteps{}-step check in which
CORAL applied \dlCoralCheckAppliedMin{}--\dlCoralCheckAppliedMax{} of \dlCoralCheckSteps{} updates
on every fold, and it fired on no run.
We hashed the analysis plan before computing any value. Appendix~\ref{app:dilute10} gives every
run's skipped updates and applied fraction.
